% GENERATED FILE. Edit canonical lessons, then run npm run generate:books.
% canonical-source-hash: fnv1a64:7a9d4485dee67154
% canonical-lessons: UR-C332-nuqsan, UR-C332-zariya, UR-C332-radd-e-amal, UR-C332-bigarna, UR-C332-phailna

\chapter{What Follows From It}
\label{ch:ur-what-follows-from-it}
\hlchaptermodality{332}

\begin{chapteropening}
\textbf{By the end of this chapter:} \emph{I can say what an event led to.}

Complete the last lesson of chapter 332: I can say what an event led to.
\end{chapteropening}

\section[\texorpdfstring{\ur{نقصان} (nuqsān)}{نقصان (nuqsān)}]{\ur{نقصان} (nuqsān) — harm; a loss}

\label{lesson:UR-C332-nuqsan}



\begin{quote}
Before the new one: say the Urdu for to be linked, then the Urdu for connected.
\end{quote}

\subsection*{You'll want to know: \ur{نقصان}}
\textbf{\ur{نقصان}} — \emph{nuqsān} — \textquotedblleft{}harm; a loss\textquotedblright{}.

\begin{quote}
\ur{بارش} \ur{سے} \ur{کسانوں} \ur{کو} \ur{نقصان} \ur{ہوا۔}
\end{quote}

\begin{quote}
\emph{bārish se kisānoṅ ko nuqsān huā.}
\end{quote}

\begin{quote}
The rain did the farmers harm.
\end{quote}

Masculine.

\begin{cousinweb}
\textbf{\ur{نقصان}} (\emph{nuqsān}) is usually traced, through Persian, to Arabic \emph{nuqṣān}, \textquotedblleft{}a decrease, a loss\textquotedblright{}.
\end{cousinweb}

\begin{grammarlens}[title={What you've built}]
One more word for giving reasons and opinions.
\end{grammarlens}

\subsection*{Your turn}
\begin{itemize}
\raggedright
  \item \emph{Say it:} \emph{nuqsān}
  \item \emph{Say it:} \emph{nuqsān}, once more
  \item \emph{Say it:} \emph{bārish se kisānoṅ ko nuqsān huā.} — The rain did the farmers harm.
  \item \emph{Recall:} say the Urdu for to be linked, then the Urdu for connected, then say \emph{nuqsān} again
\end{itemize}

\begin{tcolorbox}[breakable,colback=teal!4,colframe=teal!35!black,title={Before you move on}]
What does \ur{نقصان} mean? (\textquotedblleft{}Harm; a loss\textquotedblright{}.) Say it once more.
\end{tcolorbox}
\section[\texorpdfstring{\ur{ذریعہ} (zariyā)}{ذریعہ (zariyā)}]{\ur{ذریعہ} (zariyā) — a means; a way (of doing)}

\label{lesson:UR-C332-zariya}



\begin{quote}
Before the new one: say the Urdu for connected, then the Urdu for harm.
\end{quote}

\subsection*{You'll want to know: \ur{ذریعہ}}
\textbf{\ur{ذریعہ}} — \emph{zariyā} — \textquotedblleft{}a means; a way (of doing)\textquotedblright{}.

\begin{quote}
\ur{پیسہ} \ur{کمانے} \ur{کا} \ur{یہ} \ur{اچھا} \ur{ذریعہ} \ur{ہے۔}
\end{quote}

\begin{quote}
\emph{paisā kamāne kā yih achchhā zariyā hai.}
\end{quote}

\begin{quote}
This is a good way of earning money.
\end{quote}

Masculine.

\begin{cousinweb}
\textbf{\ur{ذریعہ}} (\emph{zariyā}) is usually traced, through Persian, to Arabic \emph{ḏarīʿa}, \textquotedblleft{}a means\textquotedblright{}.
\end{cousinweb}

\begin{grammarlens}[title={What you've built}]
One more word for giving reasons and opinions.
\end{grammarlens}

\subsection*{Your turn}
\begin{itemize}
\raggedright
  \item \emph{Say it:} \emph{zariyā}
  \item \emph{Say it:} \emph{zariyā}, once more
  \item \emph{Say it:} \emph{paisā kamāne kā yih achchhā zariyā hai.} — This is a good way of earning money.
  \item \emph{Recall:} say the Urdu for connected, then the Urdu for harm, then say \emph{zariyā} again
\end{itemize}

\begin{tcolorbox}[breakable,colback=teal!4,colframe=teal!35!black,title={Before you move on}]
What does \ur{ذریعہ} mean? (\textquotedblleft{}A means; a way (of doing)\textquotedblright{}.) Say it once more.
\end{tcolorbox}
\section[\texorpdfstring{\ur{رد} \ur{عمل} (radd-e-amal)}{رد عمل (radd-e-amal)}]{\ur{رد} \ur{عمل} (radd-e-amal) — a reaction}

\label{lesson:UR-C332-radd-e-amal}



\begin{quote}
Before the new one: say the Urdu for harm, then the Urdu for a means.
\end{quote}

\subsection*{You'll want to know: \ur{رد} \ur{عمل}}
\textbf{\ur{رد} \ur{عمل}} — \emph{radd-e-amal} — \textquotedblleft{}a reaction\textquotedblright{}.

\begin{quote}
\ur{اس} \ur{خبر} \ur{پر} \ur{لوگوں} \ur{کا} \ur{رد} \ur{عمل} \ur{کیا} \ur{تھا؟}
\end{quote}

\begin{quote}
\emph{is khabar par logoṅ kā radd-e-amal kyā thā?}
\end{quote}

\begin{quote}
How did people react to this news?
\end{quote}

Masculine.

\begin{cousinweb}
\textbf{\ur{رد} \ur{عمل}} (\emph{radd-e-amal}) is built from two loans, each usually traced, through Persian, to Arabic: Arabic \emph{radd}, \textquotedblleft{}a turning back\textquotedblright{}, and Arabic \emph{ʿamal}, \textquotedblleft{}an action\textquotedblright{}. An action sent back.
\end{cousinweb}

\begin{grammarlens}[title={What you've built}]
One more word for giving reasons and opinions.
\end{grammarlens}

\subsection*{Your turn}
\begin{itemize}
\raggedright
  \item \emph{Say it:} \emph{radd-e-amal}
  \item \emph{Say it:} \emph{radd-e-amal}, once more
  \item \emph{Say it:} \emph{is khabar par logoṅ kā radd-e-amal kyā thā?} — How did people react to this news?
  \item \emph{Recall:} say the Urdu for harm, then the Urdu for a means, then say \emph{radd-e-amal} again
\end{itemize}

\begin{tcolorbox}[breakable,colback=teal!4,colframe=teal!35!black,title={Before you move on}]
What does \ur{رد} \ur{عمل} mean? (\textquotedblleft{}A reaction\textquotedblright{}.) Say it once more.
\end{tcolorbox}
\section[\texorpdfstring{\ur{بگاڑنا} (bigāṛnā)}{بگاڑنا (bigāṛnā)}]{\ur{بگاڑنا} (bigāṛnā) — to spoil; to make worse}

\label{lesson:UR-C332-bigarna}



\begin{quote}
Before the new one: say the Urdu for a means, then the Urdu for a reaction.
\end{quote}

\subsection*{You'll want to know: \ur{بگاڑنا}}
\textbf{\ur{بگاڑنا}} — \emph{bigāṛnā} — \textquotedblleft{}to spoil; to make worse\textquotedblright{}.

\begin{quote}
\ur{بارش} \ur{نے} \ur{ہمارا} \ur{کھیل} \ur{بگاڑ} \ur{دیا۔}
\end{quote}

\begin{quote}
\emph{bārish ne hamārā khel bigāṛ diyā.}
\end{quote}

\begin{quote}
The rain spoiled our game.
\end{quote}

Past, with \ur{نے} (he or she …): \ur{اس} \ur{نے} \ur{کام} \ur{بگاڑا}, \emph{us ne kām bigāṛā}.

\begin{cousinweb}
\textbf{\ur{بگاڑنا}} (\emph{bigāṛnā}) is inherited: CDIAL lists the Hindi cousin \emph{bigāṛnā}, \textquotedblleft{}to spoil\textquotedblright{}, under Sanskrit \emph{vighaṭate}, \textquotedblleft{}breaks apart\textquotedblright{} (CDIAL 11673).
\end{cousinweb}

\begin{grammarlens}[title={What you've built}]
One more word for giving reasons and opinions.
\end{grammarlens}

\subsection*{Your turn}
\begin{itemize}
\raggedright
  \item \emph{Say it:} \emph{bigāṛnā}
  \item \emph{Say it:} \emph{bigāṛnā}, once more
  \item \emph{Say it:} \emph{bārish ne hamārā khel bigāṛ diyā.} — The rain spoiled our game.
  \item \emph{Recall:} say the Urdu for a means, then the Urdu for a reaction, then say \emph{bigāṛnā} again
\end{itemize}

\begin{tcolorbox}[breakable,colback=teal!4,colframe=teal!35!black,title={Before you move on}]
What does \ur{بگاڑنا} mean? (\textquotedblleft{}To spoil; to make worse\textquotedblright{}.) Say it once more.
\end{tcolorbox}
\section[\texorpdfstring{\ur{پھیلنا} (phailnā)}{پھیلنا (phailnā)}]{\ur{پھیلنا} (phailnā) — to spread}

\label{lesson:UR-C332-phailna}



\begin{quote}
Before the new one: say the Urdu for a reaction, then the Urdu for to spoil.
\end{quote}

\subsection*{You'll want to know: \ur{پھیلنا}}
\textbf{\ur{پھیلنا}} — \emph{phailnā} — \textquotedblleft{}to spread\textquotedblright{}.

\begin{quote}
\ur{یہ} \ur{خبر} \ur{شہر} \ur{میں} \ur{جلدی} \ur{پھیل} \ur{گئی۔}
\end{quote}

\begin{quote}
\emph{yih khabar shahr meṅ jaldī phail gaī.}
\end{quote}

\begin{quote}
The news spread quickly through the city.
\end{quote}

Past (it …): \ur{وہ} \ur{بیماری} \ur{پھیل} \ur{گئی}, \emph{voh bīmārī phail gaī}.

\begin{cousinweb}
\textbf{\ur{پھیلنا}} (\emph{phailnā}) is inherited: CDIAL lists the Hindi cousin \emph{phailnā}, \textquotedblleft{}to spread\textquotedblright{}, under Sanskrit \emph{prathita}, \textquotedblleft{}spread out\textquotedblright{} (CDIAL 8651).
\end{cousinweb}

\begin{grammarlens}[title={What you've built}]
One more word for giving reasons and opinions.
\end{grammarlens}

\subsection*{Your turn}
\begin{itemize}
\raggedright
  \item \emph{Say it:} \emph{phailnā}
  \item \emph{Say it:} \emph{phailnā}, once more
  \item \emph{Say it:} \emph{yih khabar shahr meṅ jaldī phail gaī.} — The news spread quickly through the city.
  \item \emph{Recall:} say the Urdu for a reaction, then the Urdu for to spoil, then say \emph{phailnā} again
\end{itemize}

\begin{tcolorbox}[breakable,colback=teal!4,colframe=teal!35!black,title={Before you move on}]
What does \ur{پھیلنا} mean? (\textquotedblleft{}To spread\textquotedblright{}.) Say it once more.
\end{tcolorbox}
